\section{Relasi Rekursif Tiga Suku}
\label{sec:rekurensi}

Pada bagian ini dijelaskan proses mendapatkan relasi rekursif dengan menyubstitusikan solusi Frobenius ke dalam persamaan master radial \eqref{eq:pers-radial} hingga diperoleh relasi rekursif tiga suku. Koefisien relasi rekursif yang diperoleh selanjutnya digunakan sebagai dasar dalam perhitungan numerik frekuensi mode kuasinormal.

\noindent Substitusi ansatz \eqref{eq:ansatz_frobenius} ke dalam persamaan radial hingga didapatkan relasi
rekursif bagi koefisien $a_n$. Relasi rekursif ini harus diselesaikan lebih lanjut
menggunakan metode pecahan berlanjut, yang mensyaratkan relasi rekursif berbentuk tiga suku,
\begin{equation}
    \alpha_n a_{n+1} + \beta_n a_n + \gamma_n a_{n-1} = 0,
    \label{eq:three_term}
\end{equation}
dengan koefisien $\alpha_n$, $\beta_n$, dan $\gamma_n$ sebagai fungsi yang mengandung
$\omega$. Apabila relasi rekursif yang diperoleh memiliki lebih dari tiga suku, relasi
tersebut terlebih dahulu direduksi menjadi tiga suku melalui eliminasi Gaussian
\parencite{Leaver1985, Konoplya2011}.

Bentuk pecahan berlanjut diperoleh langsung dari relasi rekursif tiga suku tersebut.
Dengan mendefinisikan rasio antara dua koefisien berurutan,
\begin{equation}
    R_n = \frac{a_n}{a_{n-1}},
\end{equation}
dan membagi Persamaan~\eqref{eq:three_term} dengan $a_n$, diperoleh hubungan rekursif
\begin{equation}
    R_n = \frac{-\gamma_n}{\beta_n + \alpha_n R_{n+1}}.
    \label{eq:rasio-rekursif}
\end{equation}
Penerapan Persamaan~\eqref{eq:rasio-rekursif} secara berulang menghasilkan bentuk pecahan
yang bersusun tanpa batas, yang dikenal sebagai pecahan berlanjut \parencite{Leaver1985}.

Deret Frobenius hanya konvergen apabila koefisien $a_n$ merupakan solusi minimal dari relasi
rekursif. Konvergensi deret pada batas luar domain ($z \to 1$) setara dengan terpenuhinya kondisi batas fisik pada batas
tersebut, dan hanya tercapai untuk nilai $\omega$ diskret tertentu yang merupakan frekuensi
\textit{quasinormal modes} \parencite{Konoplya2011}. Syarat solusi minimal terpenuhi ketika
rasio $R_1 = a_1/a_0$ dari pecahan berlanjut konsisten dengan relasi rekursif pada $n = 0$,
yaitu $\alpha_0 a_1 + \beta_0 a_0 = 0$ (karena $a_{-1}$ tidak ada). Penggabungan kedua syarat
menghasilkan kondisi pecahan berlanjut
\begin{equation}
    \beta_0 = \frac{\alpha_0 \gamma_1}{\beta_1 - \cfrac{\alpha_1 \gamma_2}{\beta_2 - \cfrac{\alpha_2 \gamma_3}{\beta_3 - \cdots}}}.
    \label{eq:continued_fraction}
\end{equation}
Persamaan~\eqref{eq:continued_fraction} bersifat transenden dan diselesaikan secara numerik
dengan memotong pecahan berlanjut pada kedalaman tertentu. Setiap akar $\omega$ yang
diperoleh berpadanan dengan satu mode pada spektrum frekuensi \textit{quasinormal modes}
\parencite{Leaver1985, Konoplya2011}.

\subsection{Schwarzschild}
\label{subsec:deriv-schw}

Ansatz Frobenius untuk ruang-waktu Schwarzschild telah diperoleh pada Persamaan~\eqref{eq:ansatz_saf}. Untuk menurunkan relasi rekursif, ansatz tersebut dituliskan kembali dalam bentuk
\begin{equation}
    \Psi(r)=\underbrace{e^{i\omega r}(r-2M)^{-2iM\omega}r^{4iM\omega}}_{W(r)}\underbrace{\sum_{n=0}^{\infty}a_nz^n}_{u(z)}.
    \label{eq:ansatz-schw-full}
\end{equation} 

Variabel deret yang digunakan adalah
\begin{equation}
    z=\frac{r-r_h}{r}=1-\frac{r_h}{r}, \qquad r=\frac{r_h}{1-z},
    \label{eq:variabel-z-schw}
\end{equation}
dengan $r_h=2M$. Variabel ini memetakan horizon peristiwa ke $z=0$ dan tak hingga spasial ke $z=1$. Dengan menggunakan fungsi metrik Schwarzschild pada Persamaan~\eqref{eq:fungsi-metrik-schwarzschild}, diperoleh
\begin{equation}
    f(r)=1-\frac{2M}{r} = z.
    \label{eq:f-sama-dengan-z}
\end{equation}
Kesamaan bentuk antara $f(r)$ dan $z$ menyederhanakan transformasi persamaan radial ke dalam variabel deret.

Turunan variabel $z$ terhadap $r$ diberikan oleh
\begin{equation}
    \frac{dz}{dr}=\frac{2M}{r^2}=\frac{(1-z)^2}{2M}, \qquad \frac{d^2z}{dr^2}=-\frac{4M}{r^3}=-\frac{(1-z)^3}{2M^2}.
    \label{eq:dzdr-schw}
\end{equation}
Dengan menerapkan aturan hasil kali dan aturan rantai pada $\Psi(r)=W(r)u(z)$, diperoleh
\begin{align}
    \frac{d\Psi}{dr} &= W\left[\frac{W'}{W}u+\frac{dz}{dr}\frac{du}{dz}\right], \notag\\[4pt]
    \frac{d^2\Psi}{dr^2} &= W\left[\frac{W''}{W}u+\left(2\frac{W'}{W}\frac{dz}{dr}+\frac{d^2z}{dr^2}\right)\frac{du}{dz}+\left(\frac{dz}{dr}\right)^2\frac{d^2u}{dz^2}\right].
    \label{eq:turunan-produk-schw}
\end{align}
Turunan faktor asimtotik yang diperlukan adalah
\begin{equation}
    \frac{W'}{W}=i\omega-\frac{2iM\omega}{r-2M}+\frac{4iM\omega}{r}, \qquad \frac{W''}{W}=\frac{d}{dr}\left(\frac{W'}{W}\right)+\left(\frac{W'}{W}\right)^2.
    \label{eq:turunan-W-schw}
\end{equation}

Persamaan~\eqref{eq:turunan-produk-schw} dan \eqref{eq:turunan-W-schw} kemudian disubstitusikan ke dalam persamaan master radial \eqref{eq:pers-radial}. Setelah seluruh persamaan dibagi dengan $W(r)$, semua fungsi yang masih bergantung pada $r$ dinyatakan dalam variabel $z$ menggunakan Persamaan~\eqref{eq:variabel-z-schw}, \eqref{eq:f-sama-dengan-z}, dan \eqref{eq:dzdr-schw}. Eliminasi faktor bersama dan penyebut menghasilkan persamaan diferensial
\begin{equation}
    A(z)\frac{d^2u}{dz^2}+B(z)\frac{du}{dz}+C(z)u=0.
    \label{eq:ode-polynomial-schw}
\end{equation}
Penjabaran aljabar secara lengkap dari tahap substitusi hingga diperolehnya Persamaan~\eqref{eq:ode-polynomial-schw} diberikan pada Lampiran~\ref{app:rekuren-schw}.

Koefisien polinomial pada Persamaan~\eqref{eq:ode-polynomial-schw} diperoleh sebagai
\begin{equation}
    \begin{aligned}
        A(z) &= z^4-2z^3+z^2, \\
        B(z) &= (3-8iM\omega)z^3+(-4+16iM\omega)z^2+(1-4iM\omega)z, \\
        C(z) &= \left(1-16M^2\omega^2-8iM\omega\right)z^2+\left[-\ell(\ell+1)-1+32M^2\omega^2+8iM\omega\right]z.
    \end{aligned}
    \label{eq:koef-schw}
\end{equation}
Deret $u(z)$ pada Persamaan~\eqref{eq:ansatz-schw-full} beserta turunannya kemudian disubstitusikan ke dalam Persamaan~\eqref{eq:ode-polynomial-schw}. Berdasarkan struktur pangkat pada Persamaan~\eqref{eq:koef-schw}, pencocokan koefisien pada setiap pangkat $z$ hanya melibatkan $a_{k+1}$, $a_k$, dan $a_{k-1}$. Dengan demikian, relasi yang diperoleh memiliki bentuk tiga suku sebagaimana Persamaan~\eqref{eq:three_term}, dengan koefisien
\begin{equation}
    \begin{aligned}
        \alpha_k &= (k+1)(k+1-4iM\omega), \\
        \beta_k &= -2k(k+1)-\ell(\ell+1)-1+32M^2\omega^2+8iM\omega(2k+1), \\
        \gamma_k &= (k-4iM\omega)^2.
    \end{aligned}
    \label{eq:koef-rekurensi-schw}
\end{equation}

\subsection{Schwarzschild Anti-de Sitter}
\label{subsec:deriv-sads}

\noindent Ansatz Frobenius untuk ruang-waktu Schwarzschild Anti-de Sitter telah diperoleh pada
Persamaan~\eqref{eq:ansatz_sads}. Untuk menurunkan relasi rekursif, langkah berikutnya adalah mensubstitusikan solusi Frobenius tersebut ke dalam persamaan master radial. Seluruh turunan terhadap koordinat $r$ kemudian ditransformasikan ke dalam variabel $z$. Selanjutnya, persamaan dikalikan dengan faktor yang sesuai untuk menghilangkan penyebut sehingga diperoleh persamaan diferensial berkoefisien polinomial,
\begin{equation}
    \widetilde{A}(z)\frac{d^2u}{dz^2} + \widetilde{B}(z)\frac{du}{dz} + \widetilde{C}(z)u = 0, \quad u(z)=\sum_{n=0}^{\infty}a_n z^n 
\end{equation}
Konvolusi deret yang muncul selama proses substitusi menghasilkan persamaan jalinan indeks untuk setiap basis, yang selanjutnya digunakan untuk memperoleh relasi rekursif.

Derajat polinomial $\widetilde{A}(z)$, $\widetilde{B}(z)$, dan $\widetilde{C}(z)$
pada kasus SAdS lebih tinggi dibandingkan Schwarzschild karena fungsi metrik $f(r)$ memuat
suku tambahan $r^2/L^2$ yang berasal dari konstanta kosmologi. Suku ini menyumbang pangkat
$r$ yang lebih tinggi pada pembilang persamaan radial, sehingga ketika seluruh persamaan
ditransformasikan ke variabel $z$, derajat polinomial koefisiennya turut meningkat.
Akibatnya, jumlah suku pada relasi rekursif yang dihasilkan pun bertambah, sebagaimana
diuraikan pada Lampiran~\ref{app:rekuren-sads}.

Pada kasus Schwarzschild Anti-de Sitter, derajat polinomial koefisien yang lebih tinggi menyebabkan relasi rekursif yang diperoleh berbentuk $N$ suku,
\begin{equation}
    \sum_{j=0}^{\min(N-1,k)} c_{j,k}^{(N)}(\omega)\, a_{k-j} = 0,
    \label{eq:rekurensi-N-suku}
\end{equation}
dengan $N>3$, sehingga tidak dapat langsung diselesaikan dengan metode pecahan berlanjut yang mensyaratkan bentuk tiga suku. Pada Persamaan~\eqref{eq:rekurensi-N-suku}, indeks $k$ menyatakan suku deret yang sedang ditinjau (sejalan dengan indeks $a_k$), sedangkan indeks $j$ menyatakan jarak suku $a_{k-j}$ terhadap $a_k$ dalam relasi tersebut. Oleh karena itu, diperlukan metode tambahan Eliminasi Gaussian untuk mereduksi relasi rekursif tersebut menjadi tiga suku.

Metode Eliminasi Gaussian yang telah disebutkan pada Subbab~\ref{sec:rekurensi} diterapkan secara bertahap untuk menghilangkan koefisien berindeks $j$ terluar, dimulai dari relasi $N$ suku hingga tersisa tiga suku \parencite{Konoplya2011}. Setiap tahap reduksi mengubah relasi $(m+1)$ suku menjadi relasi $m$ suku, dengan $m$ berkurang secara berurutan dari $N-1$ hingga $3$. Koefisien pada tahap ke-$m$ diperoleh dari koefisien pada tahap sebelumnya melalui
\begin{equation}
    c_{j,k}^{(m)}(\omega) = c_{j,k}^{(m+1)}(\omega), \qquad j=0 \ \text{atau}\ k<m,
\end{equation}
\begin{equation}
    c_{j,k}^{(m)}(\omega) = c_{j,k}^{(m+1)}(\omega)
    - \frac{c_{m,k}^{(m+1)}(\omega)\, c_{j-1,k-1}^{(m)}(\omega)}{c_{m-1,k-1}^{(m)}(\omega)},
    \qquad j\geq 1 \ \text{dan}\ k\geq m.
\end{equation}
Baris pertama berlaku untuk suku yang belum tersentuh proses eliminasi pada tahap tersebut, sedangkan baris kedua menghitung koreksi bagi suku yang terpengaruh akibat penghapusan koefisien terluar $c_{m,k}^{(m+1)}$. Iterasi ini diterapkan berulang hingga diperoleh relasi rekursif tiga suku pada $m=3$.

Hasil reduksi tersebut dapat dituliskan dalam bentuk

\begin{equation}
    \tilde{\alpha}_k a_{k+1}
    +
    \tilde{\beta}_k a_k
    +
    \tilde{\gamma}_k a_{k-1}
    =0,
\end{equation}

dengan $\tilde{\alpha}_k \equiv c_{0,k}^{(3)}$, $\tilde{\beta}_k \equiv c_{1,k}^{(3)}$, dan $\tilde{\gamma}_k \equiv c_{2,k}^{(3)}$ merupakan koefisien hasil reduksi Eliminasi Gaussian pada tahap terakhir. Berbeda dengan latar belakang Schwarzschild, relasi rekursif tiga suku ini diselesaikan secara langsung melalui rekursi maju untuk memperoleh setiap koefisien $a_k$, kemudian dijumlahkan sebagai deret; akar $\omega$ yang membuat jumlah deret tersebut lenyap pada batas anti-de Sitter, sesuai kondisi Dirichlet pada Persamaan~\eqref{eq:ansatz_sads}, merupakan frekuensi kompleks mode kuasinormal \parencite{Horowitz2000}.